Rolf pikkies recensies maakte de klik namelijk dat de klerk de apostel dames met pikkies kwalificeerde als hoeveel mannen ze tegelijk aankon. Daarop kwam eigenlijk de kliuk tussen de rentmeester en de klerk.

Hij is een matennaaier dus een homo met de rentmeester
hij was eerst De volgende grap: “Waarom praat met mij alleen over geld?” zei de tollenaar

de tollenaar verbaast zich erover dat iedereen alleen met hem over geld praat, tot iemand hem eraan herinnert wat zijn rol is (“je bent toch tollenaar?”), waarna hij daar even over nadenkt (“en hij nadacht”).
De grap draait om het stereotype dat iemand zijn eigen imago meestal niet helemaal doorziet, totdat iemand het simpel benoemt.

De klerk koos een euvel
en sloeg dan een deuvel

in zijn zelfgeboord gat,
zodat men hem aanbad

voor zijn vernuftigheid.
Pas later kreeg men spijt.

De rentmeester hapte,
omdat hij goed snapte

dit politiek machtsspel.
Hij waardeerde het wel.

Hij vond het kerkelijk
en woest aantrekkelijk.

De klerk voelde snel aan
dat hij naar bed moest gaan

ook met de rentmeester.
Zo werd hij de meester.

“Hij is net als een fret”,
zei de ‘Profeet’ zonet,

tegen de kerkenraad
over de ‘Klerk‘’s verraad.

“Hij slipt er weer doorheen
en glipt er weer gemeen

heel gladjes tussenuit.”
De ‘Rentmeester’ gaf luid

een kattige sisklank.
“Dank voor de stank voor dank

voor mijn waarschuwingen.
U kunt het verdringen

dat hij over de raad
achter uw rug om kwaad

en denigrerend spreekt.
Het is uw dolk die steekt

etterend in uw rug.
En u reageert stug

en ziet niet zijn motief
als gemaskeerde dief.

De apostelen zijn
voor hem een feestfestijn.

Hij vroeg me al kwijlend
op hun schoonheid geilend:

“Mocht jij ze eerst keuren
in al hun haarkleuren

onder de regenboog?”
Het behoeft geen betoog

dat hij mijn leerlingen
alleen ziet als dingen

voor zijn verdorven lust.
Kijk eens hoe hij ze kust

en ongevraagd betast.
Dat is toch ongepast?

De kracht van metalen
deed hun naam bepalen.

En God heeft die zowaar
aan de kleur van hun haar

weten te verbinden.
Niet om te verblinden!

Ze zijn verkondigers
alsook verdedigers

van Pandemonium.
Geef hen het podium

voor hun levensstijlen,
niet om op te geilen.

De 'Klerk' is ook corrupt.
Hij belooft u abrupt

financiële support.
Waar het echter aan schort

is dat hij een groot deel
eist als eigen aandeel.

De ‘Klerk’ hamstert en steelt
wat je ook met hem deelt.

Hij dekt zich steeds listig.
De ‘Klerk’ presenteert zich

als een heilig boontje.
Zonder een juist loontje

als hij zich weer misdraagt.
En dan steen en been klaagt

over de gebreken,
die zogenaamd bleken,

bij een ander persoon.
Uw dufheid is zijn loon.”

De ‘Klerk’’s juridisme
met whataboutisme:

“Wie denk je dat je bent?”
“Iemand die zich niet kent,

maar elke dag nog leert
wat hem pleziert of deert.”

waarmee hij, rechtsgeleerd, zich diskwalificeert

volgens de ijkmaten
van de technocraten.

De ‘Klerk’ werd dus geroemd
en uit achting benoemd

als ‘Rechter’ van de kerk.
Een sterk staaltje praatwerk.

Rechter oneerlijk gezegd
